\setcounter{chapter}{23}
\renewcommand{\thetheorem}{24.a.\arabic{theorem}}
\setcounter{theorem}{0}

\chapter{Spectral Theorems}

\begin{ainote}
The professor left a note to \qt{explain the name} of the Spectral Theorems. In linear algebra, the \qt{spectrum} of a linear operator is its set of eigenvalues. The Spectral Theorems provide conditions under which an operator can be completely described by its spectrum—specifically, when there exists a basis for the entire space consisting solely of the operator's eigenvectors (making the operator diagonalizable).
\end{ainote}

% Def. (24.a.1)
\begin{definition}
\label{def:orthogonally_diagonalizable}
Let $V$ be an inner product space and $T \colon V \to V$ an endomorphism. We say that $T$ is \newterm{orthogonally diagonalizable} if there exists an orthonormal basis for $V$ consisting of eigenvectors of $T$.
\end{definition}

% Lemma. (24.a.2)
\begin{lemma}
\label{lem:orth_diag_implies_commute}
If $T$ is orthogonally diagonalizable, then $T \circ T^* = T^* \circ T$.
\end{lemma}

\begin{proof}
Let $\mathcal{B}$ be an orthonormal basis for $V$ such that all the vectors in $\mathcal{B}$ are eigenvectors of $T$. This implies that the representation matrix $[T]_{\mathcal{B}}^{\mathcal{B}}$ is a diagonal matrix. 

Since $\mathcal{B}$ is an orthonormal basis, we also know that the matrix of the adjoint map is the conjugate transpose of the matrix of the map, meaning $[T^*]_{\mathcal{B}}^{\mathcal{B}} = \left([T]_{\mathcal{B}}^{\mathcal{B}}\right)^*$. Since the conjugate transpose of a diagonal matrix is still a diagonal matrix, $[T^*]_{\mathcal{B}}^{\mathcal{B}}$ is also a diagonal matrix. 

If $M$ and $N$ are diagonal matrices, then they must commute, meaning $M \cdot N = N \cdot M$. Therefore, we have:
\[ [T]_{\mathcal{B}}^{\mathcal{B}} \cdot [T^*]_{\mathcal{B}}^{\mathcal{B}} = [T^*]_{\mathcal{B}}^{\mathcal{B}} \cdot [T]_{\mathcal{B}}^{\mathcal{B}}. \]
By the properties of representation matrices, this implies:
\[ [T \circ T^*]_{\mathcal{B}}^{\mathcal{B}} = [T^* \circ T]_{\mathcal{B}}^{\mathcal{B}} \implies T \circ T^* = T^* \circ T. \]
\end{proof}

% Def. (24.a.3)
\begin{definition}
\label{def:normal_endomorphism}
Let $V$ be an inner product space. An endomorphism $T \colon V \to V$ is called \newterm{normal} if $T \circ T^* = T^* \circ T$. 

Similarly, a matrix $A \in \M_{n \times n}(K)$ (where $K = \mathbb{R}$ or $\mathbb{C}$) is called \newterm{normal} if $A \cdot \overline{A}^\top = \overline{A}^\top \cdot A$. (Note that $A$ is a normal matrix if and only if the map $T_A \colon K^n \to K^n$ is a normal endomorphism, provided we endow $K^n$ with its standard inner product).
\end{definition}

We have seen in the previous lemma that if $T \colon V \to V$ is orthogonally diagonalizable, then $T$ is normal. 

\begin{question}
Is the converse statement true?
\end{question}

Over $K = \mathbb{R}$, the answer is NO! 
For example, consider the matrix $A = \left(\begin{smallmatrix} 0 & 1 \\ -1 & 0 \end{smallmatrix}\right)$. The linear map $T_A \colon \mathbb{R}^2 \to \mathbb{R}^2$ is a clockwise rotation by $90^\circ$, so there are no eigenvalues (in $\mathbb{R}$) and no eigenvectors. Thus, it cannot be diagonalizable. 
However, $\overline{A}^\top = A^\top = \left(\begin{smallmatrix} 0 & -1 \\ 1 & 0 \end{smallmatrix}\right)$, and we can easily check that $A \cdot A^\top = A^\top \cdot A = I_2$, meaning $A$ is indeed normal.

What about over $\mathbb{C}$?

% Thm. (24.a.4)
\begin{theorem}[Spectral Theorem over $\mathbb{C}$]
\label{thm:spectral_thm_c}
Let $V$ be a finite-di\-men\-sional inner product space over $\mathbb{C}$, and let $T \colon V \to V$ be a linear map. Then $T$ is orthogonally diagonalizable if and only if $T$ is normal.
\end{theorem}

For the proof, we need a lemma.

% Lemma. (24.a.5)
\begin{lemma}[Properties of Normal Endomorphisms]
\label{lem:properties_normal_endo}
Let $V$ be a finite-di\-men\-sional inner product space over $K$ (where $K = \mathbb{R}$ or $\mathbb{C}$). Let $T \colon V \to V$ be a normal endomorphism. Then:
\begin{enumerate}
    \item \textbf{(a)} For all $v \in V$, $\|Tv\| = \|T^*v\|$.
    \item \textbf{(b)} For all $\lambda \in K$, the endomorphism $T - \lambda \id_V$ is also normal.
    \item \textbf{(c)} If $v$ is an eigenvector of $T$ with eigenvalue $\lambda$, then $v$ is also an eigenvector of $T^*$ with eigenvalue $\overline{\lambda}$.
    \item \textbf{(d)} Let $v_1, v_2$ be eigenvectors of $T$ with eigenvalues $\lambda_1, \lambda_2$ respectively. If $\lambda_1 \neq \lambda_2$, then $\langle v_1, v_2 \rangle = 0$ (i.e., $v_1 \perp v_2$).
\end{enumerate}
\end{lemma}

\begin{proof}
\leavevmode
\begin{description}
    \item[Proof of (a):] We compute the squared norm:
    \[ \|Tv\|^2 = \langle Tv, Tv \rangle = \langle v, T^*Tv \rangle = \langle v, TT^*v \rangle = \langle T^*v, T^*v \rangle = \|T^*v\|^2, \]
    where we used the definition of the adjoint and the fact that $T$ is normal ($T \circ T^* = T^* \circ T$).
    
    \item[Proof of (b):] We expand the compositions:
    \begin{align*}
        (T - \lambda \id_V) \circ (T - \lambda \id_V)^* &= (T - \lambda \id_V) \circ (T^* - \overline{\lambda} \id_V) \\
        &= T \circ T^* - \overline{\lambda} T - \lambda T^* + \lambda \overline{\lambda} \id_V. \quad (*)
    \end{align*}
    On the other hand:
    \begin{align*}
        (T - \lambda \id_V)^* \circ (T - \lambda \id_V) &= (T^* - \overline{\lambda} \id_V) \circ (T - \lambda \id_V) \\
        &= T^* \circ T - \lambda T^* - \overline{\lambda} T + \overline{\lambda} \lambda \id_V. \quad (**)
    \end{align*}
    Because $T \circ T^* = T^* \circ T$, the expressions $(*)$ and $(**)$ are equal.
    
    \item[Proof of (c):] If $v$ is an eigenvector of $T$ with eigenvalue $\lambda$, then $(T - \lambda \id_V)v = 0$. By statement \textbf{(b)}, $T - \lambda \id_V$ is normal, hence by statement \textbf{(a)}, we have:
    \[ \|(T^* - \overline{\lambda} \id_V)v\| = \|(T - \lambda \id_V)v\| = \|0\| = 0. \]
    This implies that $(T^* - \overline{\lambda} \id_V)v = 0$, and therefore $T^*v = \overline{\lambda}v$.
    
    \item[Proof of (d):] We compute the inner product in two ways:
    \[ \lambda_1 \langle v_1, v_2 \rangle = \langle \lambda_1 v_1, v_2 \rangle = \langle Tv_1, v_2 \rangle = \langle v_1, T^*v_2 \rangle. \]
    By statement \textbf{(c)}, $T^*v_2 = \overline{\lambda_2}v_2$. Thus:
    \[ \langle v_1, T^*v_2 \rangle = \langle v_1, \overline{\lambda_2} v_2 \rangle = \lambda_2 \langle v_1, v_2 \rangle. \]
    Equating the two results gives $\lambda_1 \langle v_1, v_2 \rangle = \lambda_2 \langle v_1, v_2 \rangle$, which implies $(\lambda_1 - \lambda_2) \langle v_1, v_2 \rangle = 0$. 
    Since $\lambda_1 \neq \lambda_2$, we must have $\langle v_1, v_2 \rangle = 0$.
\end{description}
\end{proof}

\begin{proof}[Proof of the Spectral Theorem over $\mathbb{C}$ (\cref{thm:spectral_thm_c})]
We only need to prove the direction \qt{$T$ is normal $\implies T$ is orthogonally diagonalizable}, as the reverse direction was already established in \cref{lem:orth_diag_implies_commute}.

By a previous theorem, $T$ is trigonalizable (here we explicitly use the fact that $K = \mathbb{C}$ and the Fundamental Theorem of Algebra). That is, there exists a basis $\mathcal{C}$ for $V$ such that $[T]_{\mathcal{C}}^{\mathcal{C}}$ is an upper triangular matrix. 
By applying the Gram-Schmidt orthogonalization process to the basis $\mathcal{C}$, we obtain a new basis $\mathcal{B} = (v_1, \dots, v_n)$ which is orthonormal, and such that $[T]_{\mathcal{B}}^{\mathcal{B}}$ is still upper triangular. (This is left as an exercise to check. The main point is that if $\mathcal{C} = (u_1, \dots, u_n)$, then Gram-Schmidt gives an orthonormal basis $\mathcal{B} = (v_1, \dots, v_n)$ with $\Sp(v_1, \dots, v_k) = \Sp(u_1, \dots, u_k)$ for all $1 \leq k \leq n$).

We claim that $[T]_{\mathcal{B}}^{\mathcal{B}}$ is actually diagonal. To see this, let us write the matrix explicitly as $[T]_{\mathcal{B}}^{\mathcal{B}} = (a_{ij})_{1 \leq i \leq n, 1 \leq j \leq n}$:
\[ [T]_{\mathcal{B}}^{\mathcal{B}} = \begin{pmatrix} a_{11} & a_{12} & \dots & a_{1n} \\ 0 & a_{22} & \dots & a_{2n} \\ \vdots & \vdots & \ddots & \vdots \\ 0 & 0 & \dots & a_{nn} \end{pmatrix}. \]
Since $\mathcal{B}$ is orthonormal, the norm of $Tv_1$ is simply the norm of the first column of the matrix, which means $\|Tv_1\|^2 = |a_{11}|^2$. 
Since $[T^*]_{\mathcal{B}}^{\mathcal{B}} = \left([T]_{\mathcal{B}}^{\mathcal{B}}\right)^*$, the first column of $[T^*]_{\mathcal{B}}^{\mathcal{B}}$ is the conjugate transpose of the first row of $[T]_{\mathcal{B}}^{\mathcal{B}}$. Therefore, we also have:
\[ \|T^*v_1\|^2 = \sum_{k=1}^n |a_{1k}|^2 = |a_{11}|^2 + |a_{12}|^2 + \dots + |a_{1n}|^2. \]
As $T$ is normal, \cref{lem:properties_normal_endo} part \textbf{(a)} tells us that $\|Tv_1\| = \|T^*v_1\|$. This implies:
\[ |a_{11}|^2 = |a_{11}|^2 + |a_{12}|^2 + \dots + |a_{1n}|^2, \]
hence $a_{12} = \dots = a_{1n} = 0$. 
This shows that the matrix now looks like:
\[ [T]_{\mathcal{B}}^{\mathcal{B}} = \begin{pmatrix} a_{11} & 0 & \dots & 0 \\ 0 & a_{22} & \dots & a_{2n} \\ \vdots & 0 & \ddots & \vdots \\ 0 & 0 & \dots & a_{nn} \end{pmatrix}. \]
Now, consider $v_2$. We have $\|Tv_2\|^2 = |a_{22}|^2$ and $\|T^*v_2\|^2 = |a_{22}|^2 + |a_{23}|^2 + \dots + |a_{2n}|^2$. 
As $\|Tv_2\| = \|T^*v_2\|$, we similarly obtain $a_{23} = \dots = a_{2n} = 0$. 

We can continue this process by induction, and show that for all $1 \leq k < n$ we have $a_{k, k+1} = \dots = a_{k,n} = 0$. Hence, $[T]_{\mathcal{B}}^{\mathcal{B}}$ is diagonal. Since $\mathcal{B}$ is an orthonormal basis that diagonalizes $T$, it consists of eigenvectors, proving that $T$ is orthogonally diagonalizable.
\end{proof}

\begin{remark}
In the proof above, we used $K = \mathbb{C}$ \underline{only} at one place: when we stated that $T$ is trigonalizable. Therefore, we immediately obtain the following corollary for the real case.
\end{remark}

% Cor. (24.a.6)
\begin{corollary}
\label{cor:spectral_thm_real_trigonalizable}
Let $V$ be a finite-di\-men\-sional inner product space over $\mathbb{R}$, and $T \colon V \to V$ an endomorphism which is normal. If $T$ is trigonalizable (which is true if and only if the characteristic polynomial $P_T(x)$ factors as a product of linear factors over $\mathbb{R}$), then $T$ is orthogonally diagonalizable.
\end{corollary}

\subsection{The Spectral Theorem over \texorpdfstring{$\mathbb{R}$}{R}}

We've seen that over $\mathbb{R}$, the statement \qt{$T \in \End(V)$ is normal $\implies T$ is orthogonally diagonalizable} does not hold in general. We need a stronger condition.

% Lemma. (24.a.7)
\begin{lemma}
\label{lem:real_orth_diag_is_self_adj}
Let $V$ be a finite-di\-men\-sional inner product space over $\mathbb{R}$, and $T \in \End(V)$ which is orthogonally diagonalizable. Then $T^* = T$.
\end{lemma}

\begin{proof}
Let $\mathcal{B}$ be an orthonormal basis such that $[T]_{\mathcal{B}}^{\mathcal{B}}$ is a diagonal matrix. Since we are over $\mathbb{R}$, the conjugate transpose is simply the standard transpose. Then:
\[ [T^*]_{\mathcal{B}}^{\mathcal{B}} = \left([T]_{\mathcal{B}}^{\mathcal{B}}\right)^\top = [T]_{\mathcal{B}}^{\mathcal{B}}. \]
(The last equality holds because $[T]_{\mathcal{B}}^{\mathcal{B}}$ is diagonal). Because their representation matrices with respect to the same basis are equal, it follows that $T^* = T$.
\end{proof}

% Def. (24.a.8)
\begin{definition}
\label{def:self_adjoint}
Let $V$ be an inner product space over $K$ (where $K = \mathbb{R}$ or $\mathbb{C}$), and $T \in \End(V)$. $T$ is called \newterm{self adjoint} if $T^* = T$. 

Similarly, a matrix $A \in \M_{n \times n}(K)$ is called \newterm{self adjoint} if $A^* = A$. (Note that over $K = \mathbb{R}$, a self adjoint matrix is precisely a symmetric matrix, $A^\top = A$).
\end{definition}

\begin{exercise}
Show that $T \colon V \to V$ is self adjoint if and only if for every orthonormal basis $\mathcal{B}$ of $V$, the matrix $[T]_{\mathcal{B}}^{\mathcal{B}}$ is self adjoint.
\end{exercise}

% Thm. (24.a.9)
\begin{theorem}[Spectral Theorem over $\mathbb{R}$]
\label{thm:spectral_thm_r}
Let $V$ be a Euclidean space (i.e., an inner product space over $\mathbb{R}$) of finite dimension. Then $T$ is orthogonally diagonalizable if and only if $T$ is self adjoint.
\end{theorem}

For the proof, we need the following lemma:

% Lemma. (24.a.10)
\begin{lemma}
\label{lem:self_adj_properties}
Let $V$ be an inner product space over $K$ (where $K = \mathbb{R}$ or $\mathbb{C}$), and let $T \in \End(V)$ be self adjoint. Then:
\begin{enumerate}
    \item \textbf{(a)} All eigenvalues of $T$ are real.
    \item \textbf{(b)} The characteristic polynomial $P_T(x)$ factors into a product of linear factors over $K$.
\end{enumerate}
\end{lemma}

\begin{proof}
\leavevmode
\begin{description}
    \item[Proof of (a):] If $K = \mathbb{R}$, there is nothing to show. So assume $K = \mathbb{C}$. Let $\lambda \in \mathbb{C}$ be an eigenvalue and $v \in V$ a corresponding non-zero eigenvector. We've seen in \cref{lem:properties_normal_endo} \textbf{(c)} that $v$ is also an eigenvector of $T^*$ with eigenvalue $\overline{\lambda}$. But $T^* = T$. This implies:
    \[ \lambda v = Tv = T^*v = \overline{\lambda} v. \]
    Since $v \neq 0$, we must have $\lambda = \overline{\lambda}$, which implies $\lambda \in \mathbb{R}$.
    
    \item[Proof of (b):] For $K = \mathbb{C}$, this follows immediately from the Fundamental Theorem of Algebra. So assume $K = \mathbb{R}$. Let $\mathcal{B}$ be an orthonormal basis for $V$ and consider the matrix $A := [T]_{\mathcal{B}}^{\mathcal{B}} \in \M_{n \times n}(\mathbb{R})$. The statement would follow if we show that the characteristic polynomial $P_A(x)$ factors as a product of linear terms with real coefficients. 
    
    Now, $T^* = T$, hence $A = [T]_{\mathcal{B}}^{\mathcal{B}} = [T^*]_{\mathcal{B}}^{\mathcal{B}} = \left([T]_{\mathcal{B}}^{\mathcal{B}}\right)^* = A^* = A^\top$. 
    Consider $\mathbb{C}^n$ endowed with its standard inner product, and view $A$ as a matrix in $\M_{n \times n}(\mathbb{C})$. Let $T_A \colon \mathbb{C}^n \to \mathbb{C}^n$ be the corresponding linear map. The standard basis $\mathcal{E}_n = (e_1, \dots, e_n)$ is orthonormal, and $[T_A]_{\mathcal{E}_n}^{\mathcal{E}_n} = A$. 
    
    We know that $P_{T_A}(x) = P_A(x)$. When viewed as a polynomial over $\mathbb{C}$ (i.e., in $\mathbb{C}[x]$), it factors as a product of linear terms:
    \begin{equation}
        P_A(x) = \prod_{k=1}^n (x - \lambda_k) \in \mathbb{C}[x], \quad (*)
    \end{equation}
    where $\lambda_1, \dots, \lambda_n$ are the eigenvalues of $T_A$. But $T_A$ is self adjoint (even when viewed over $\mathbb{C}$ since $A^\top = A = A^*$), so by part \textbf{(a)}, all the eigenvalues $\lambda_1, \dots, \lambda_n \in \mathbb{R}$. Therefore, the factorization $(*)$ actually holds in $\mathbb{R}[x]$.
\end{description}
\end{proof}

\begin{proof}[Proof of the Spectral Theorem over $\mathbb{R}$ (\cref{thm:spectral_thm_r})]
The forward direction was proven in \cref{lem:real_orth_diag_is_self_adj}. For the reverse direction, if $T$ is self adjoint, then by \cref{lem:self_adj_properties}, its characteristic polynomial splits over $\mathbb{R}$. This implies $T$ is trigonalizable over $\mathbb{R}$. Furthermore, since $T = T^*$, $T$ trivially commutes with its adjoint, making it normal. By \cref{cor:spectral_thm_real_trigonalizable}, a normal map that is trigonalizable over $\mathbb{R}$ is orthogonally diagonalizable.
\end{proof}

\subsection{Spectral Theorems for Matrices}

Let $K = \mathbb{R}$ or $\mathbb{C}$. Recall that a matrix $A \in \M_{n \times n}(K)$ is called normal if $A \cdot A^* = A^* \cdot A$.

\begin{example}
\leavevmode
\begin{enumerate}
    \item For $K = \mathbb{R}$: Symmetric matrices and orthogonal matrices $\operatorname{O}(n)$ are normal.
    \item For $K = \mathbb{C}$: Self adjoint matrices and unitary matrices $\operatorname{U}(n)$ are normal.
\end{enumerate}
\end{example}

\begin{remark}
There exist symmetric matrices in $\M_{n \times n}(\mathbb{C})$ that are not normal. For example, let $A = \left(\begin{smallmatrix} i & i \\ i & 1 \end{smallmatrix}\right)$. We have $A^\top = A$, so it is symmetric. However, $A^* = \left(\begin{smallmatrix} -i & -i \\ -i & 1 \end{smallmatrix}\right)$. A quick calculation shows that $A \cdot A^* \neq A^* \cdot A$, so $A$ is not normal.
\end{remark}

% Thm. (24.a.11)
\begin{theorem}[Matrix Spectral Theorems]
\label{thm:spectral_thm_matrices}
\leavevmode
\begin{enumerate}[label=\textbf{(\alph*)}]
    \item If $A \in \M_{n \times n}(\mathbb{C})$ is normal, then there exists a unitary matrix $U \in \operatorname{U}(n)$ such that $U^{-1} A U$ is diagonal. (Note that $U^{-1} = U^*$).
    \item If $A \in \M_{n \times n}(\mathbb{R})$ is a symmetric matrix, then there exists an orthogonal matrix $O \in \operatorname{O}(n)$ such that $O^{-1} A O$ is diagonal. (Note that $O^{-1} = O^\top$).
\end{enumerate}
\end{theorem}

\begin{proof}[Proof Sketch]
The proof is left as an exercise, based on the following key points translating between maps and matrices:
\begin{enumerate}
    \item[(a)] $A \in \M_{n \times n}(\mathbb{C})$ is normal if and only if $T_A \colon \mathbb{C}^n \to \mathbb{C}^n$ is normal (when we endow $\mathbb{C}^n$ with its standard inner product).
    \item[(b)] If $\mathcal{B}$ is an orthonormal basis for $\mathbb{C}^n$, then the transition matrix $U := [\id_{\mathbb{C}^n}]_{\mathcal{E}_n}^{\mathcal{B}}$ is unitary.
    \item[(c)] $A \in \M_{n \times n}(\mathbb{R})$ is symmetric if and only if $T_A \colon \mathbb{R}^n \to \mathbb{R}^n$ is self adjoint, when we endow $\mathbb{R}^n$ with its standard inner product.
    \item[(d)] If $\mathcal{B}$ is an orthonormal basis for $\mathbb{R}^n$, then the transition matrix $O := [\id_{\mathbb{R}^n}]_{\mathcal{E}_n}^{\mathcal{B}}$ is orthogonal.
\end{enumerate}
\end{proof}

\section*{Exercise Solutions}
\addcontentsline{toc}{section}{Exercise Solutions}

\begin{exercisesolution}[Proof of Gram-Schmidt Preserving Upper Triangularity (on \cpageref{thm:spectral_thm_c})]~\\
Let $\mathcal{C} = (u_1, \dots, u_n)$ be a basis for $V$ such that the representation matrix $[T]_{\mathcal{C}}^{\mathcal{C}}$ is upper triangular. This means that for any $1 \leq k \leq n$, the vector $T(u_k)$ is a linear combination of $u_1, \dots, u_k$. In other words, $T(u_k) \in \Sp(u_1, \dots, u_k)$. Consequently, the subspace $W_k := \Sp(u_1, \dots, u_k)$ is $T$-invariant, meaning $T(W_k) \subseteq W_k$.

Applying the Gram-Schmidt orthogonalization process to the basis $\mathcal{C}$ produces an orthonormal basis $\mathcal{B} = (v_1, \dots, v_n)$ with the property that for every $1 \leq k \leq n$, the spans are preserved: $\Sp(v_1, \dots, v_k) = \Sp(u_1, \dots, u_k) = W_k$. 

Now, consider the matrix $[T]_{\mathcal{B}}^{\mathcal{B}}$. The $k$-th column of this matrix consists of the coordinates of $T(v_k)$ with respect to the basis $\mathcal{B}$. Since $v_k \in W_k$ and $W_k$ is $T$-invariant, we must have $T(v_k) \in W_k = \Sp(v_1, \dots, v_k)$. This means $T(v_k)$ can be written strictly as a linear combination of $v_1, \dots, v_k$, meaning all coefficients for vectors $v_j$ with $j > k$ are exactly zero. Therefore, $[T]_{\mathcal{B}}^{\mathcal{B}}$ is an upper triangular matrix.
\end{exercisesolution}

\begin{exercisesolution}[Proof of Matrix Self Adjointness Equivalence (on \cpageref{def:self_adjoint})] ~\\
Let $T \in \End(V)$.
First, suppose $T$ is self adjoint, so $T = T^*$. Let $\mathcal{B}$ be an arbitrary orthonormal basis for $V$. By the properties of representation matrices for adjoint maps in orthonormal bases, we have $[T^*]_{\mathcal{B}}^{\mathcal{B}} = \left([T]_{\mathcal{B}}^{\mathcal{B}}\right)^*$. Since $T = T^*$, we substitute to get $[T]_{\mathcal{B}}^{\mathcal{B}} = \left([T]_{\mathcal{B}}^{\mathcal{B}}\right)^*$, which means the matrix $[T]_{\mathcal{B}}^{\mathcal{B}}$ is self adjoint.

Conversely, suppose that for any orthonormal basis $\mathcal{B}$ of $V$, the matrix $[T]_{\mathcal{B}}^{\mathcal{B}}$ is self adjoint, i.e., $[T]_{\mathcal{B}}^{\mathcal{B}} = \left([T]_{\mathcal{B}}^{\mathcal{B}}\right)^*$. Again, using the identity $[T^*]_{\mathcal{B}}^{\mathcal{B}} = \left([T]_{\mathcal{B}}^{\mathcal{B}}\right)^*$, we see that $[T]_{\mathcal{B}}^{\mathcal{B}} = [T^*]_{\mathcal{B}}^{\mathcal{B}}$. Because the map sending an endomorphism to its representation matrix with respect to a fixed basis is an isomorphism, this equality of matrices implies the equality of the maps, hence $T = T^*$, meaning $T$ is self adjoint.
\end{exercisesolution}

\begin{exercisesolution}[Proof of \cref{thm:spectral_thm_matrices}]
\leavevmode
\begin{description}
    \item[Proof of (a):] Let $A \in \M_{n \times n}(\mathbb{C})$ be a normal matrix. By point (a), the linear map $T_A \colon \mathbb{C}^n \to \mathbb{C}^n$ is a normal endomorphism with respect to the standard inner product. By the Spectral Theorem over $\mathbb{C}$ (\cref{thm:spectral_thm_c}), $T_A$ is orthogonally diagonalizable. Thus, there exists an orthonormal basis $\mathcal{B} = (v_1, \dots, v_n)$ for $\mathbb{C}^n$ consisting of eigenvectors of $T_A$. 
    
    The representation matrix of $T_A$ with respect to $\mathcal{B}$, which is $[T_A]_{\mathcal{B}}^{\mathcal{B}}$, is a diagonal matrix $D$. Let $\mathcal{E}_n$ be the standard basis for $\mathbb{C}^n$. The transition matrix from $\mathcal{B}$ to $\mathcal{E}_n$ is $U := [\id_{\mathbb{C}^n}]_{\mathcal{E}_n}^{\mathcal{B}}$, whose columns are precisely the vectors $v_1, \dots, v_n$. Since $\mathcal{B}$ is an orthonormal basis, point (b) guarantees that $U$ is a unitary matrix. 
    
    Using the change of basis formula, we have:
    \[ D = [T_A]_{\mathcal{B}}^{\mathcal{B}} = [\id_{\mathbb{C}^n}]_{\mathcal{B}}^{\mathcal{E}_n} \cdot [T_A]_{\mathcal{E}_n}^{\mathcal{E}_n} \cdot [\id_{\mathbb{C}^n}]_{\mathcal{E}_n}^{\mathcal{B}} = U^{-1} \cdot A \cdot U. \]
    Thus, $U^{-1} A U$ is diagonal.
    
    \item[Proof of (b):] Let $A \in \M_{n \times n}(\mathbb{R})$ be a symmetric matrix. By point (c), the linear map $T_A \colon \mathbb{R}^n \to \mathbb{R}^n$ is self adjoint. By the Spectral Theorem over $\mathbb{R}$ (\cref{thm:spectral_thm_r}), $T_A$ is orthogonally diagonalizable. Thus, there exists an orthonormal basis $\mathcal{B}$ for $\mathbb{R}^n$ consisting of eigenvectors of $T_A$, making $[T_A]_{\mathcal{B}}^{\mathcal{B}}$ a diagonal matrix $D$. 
    
    Let $O := [\id_{\mathbb{R}^n}]_{\mathcal{E}_n}^{\mathcal{B}}$ be the transition matrix. Since $\mathcal{B}$ is an orthonormal basis in $\mathbb{R}^n$, point (d) guarantees that $O$ is an orthogonal matrix. 
    By the change of basis formula:
    \[ D = [T_A]_{\mathcal{B}}^{\mathcal{B}} = [\id_{\mathbb{R}^n}]_{\mathcal{B}}^{\mathcal{E}_n} \cdot [T_A]_{\mathcal{E}_n}^{\mathcal{E}_n} \cdot [\id_{\mathbb{R}^n}]_{\mathcal{E}_n}^{\mathcal{B}} = O^{-1} \cdot A \cdot O. \]
    Thus, $O^{-1} A O$ is diagonal.
\end{description}
\end{exercisesolution}